\chapter{El orden político}
\label{ch:polity}

\section{El problema es la irrevocabilidad, no la política pública}

El capítulo~\ref{ch:ledgerch} terminó en un hallazgo, y este capítulo es su
consecuencia.

Casi todas las medidas que propone este libro ya se han intentado en Cuba.
Restaurantes privados: intentados, 1993, funcionaron, restringidos en 2017.
Mercados campesinos: intentados, 1980, funcionaron, abolidos en 1986, intentados de
nuevo en 1994. Dólares: legalizados en 1993, retirados en 2004. Tierra en
usufructo: 2008 y 2012. Empresas privadas con personalidad jurídica: 2021, y para
2024 con topes de precio y culpadas de la inflación. Al Estado cubano no le han
faltado buenas ideas. Ha sido incapaz de conservar ninguna más allá de la
emergencia que la produjo.

La consecuencia no es solo que las reformas se reviertan. Es que rinden poco
\emph{mientras están vigentes}, porque todo el mundo ve que son revocables. Un
cubano que decide si invierte los ahorros de tres años en equipamiento está
apostando a la durabilidad de un decreto, y el historial de esa apuesta es malo.
Así que no se compromete capital, los negocios se quedan pequeños y líquidos, nadie
forma a nadie, nadie construye nada que no pueda moverse ---y la mediocridad
resultante se presenta después como prueba de que el sector privado no funciona, lo
que justifica la siguiente restricción---.

Este es un problema resuelto en la economía institucional y la solución tiene nombre
\citep{northweingast1989}. Lo que un Estado necesita, cuando quiere que la gente
invierta detrás de sus promesas, es un mecanismo con el que pueda \emph{inhabilitar
su propia discrecionalidad futura} de un modo que los de fuera puedan verificar. No
una promesa mejor: una promesa que no pueda romper barato. Todo lo de este capítulo
es un instrumento de esa clase, y todo lo demás de la Parte III depende de que
existan al menos algunos de ellos. Un capítulo de excelente política energética en
un país donde la política es revocable a voluntad es un capítulo sobre lo que pudo
haber sido.

\begin{design}{El compromiso antes que el contenido}
\label{d:commit}
\begin{rules}
\rl{1}{Orden} Ninguna medida liberalizadora de la Parte~III se adopta sin al menos
un dispositivo de compromiso adherido: un anclaje constitucional, un derecho
justiciable, una obligación por tratado, un regulador independiente con mandato y
presupuesto asegurados, o una regla publicada con disparador automático. Una medida
adoptada solo por decreto la entiende todo el mundo, correctamente, como temporal.

\rl{2}{Derechos, no listas} La actividad económica se permite por derecho, sujeta a
prohibiciones publicadas, y nunca por lista enumerada de categorías autorizadas. La
lista es lo que hizo del sector privado cubano un conjunto de oficios sin ningún
desarrollador de software dentro, y es además el instrumento con el que se encoge
al sector sin derogar ninguna ley.

\rl{3}{Asimetría} Los instrumentos se diseñan de modo que ampliar una libertad sea
administrativamente barato y retirarla sea caro ---exigiendo ley, publicación,
plazos de aviso e indemnización--- y no al revés, que es el arreglo actual.

\rl{4}{Legitimación} Todo derecho creado en la Parte~III es exigible por su titular,
ante un tribunal, contra el Estado. Un derecho por el que nadie puede demandar es
una política pública, y las políticas públicas son precisamente aquello en lo que
este capítulo trata de no confiar.
\end{rules}
\end{design}

\section{La secuencia: la ley antes que las urnas}

La decisión más consecuente de este capítulo es el orden, y el libro toma una
posición que no complacerá a todos.

\emph{Estado de derecho primero; elecciones nacionales competitivas después, y con
calendario publicado.}

El razonamiento es empírico y no de principio. Las transiciones poscomunistas que
salieron bien ---Estonia, Polonia, Eslovenia, las tierras checas--- construyeron
tribunales, registros de la propiedad, autoridades de competencia y bancos centrales
independientes junto con su apertura electoral o antes de ella. Las que salieron mal
celebraron elecciones dentro de un vacío institucional, donde el ganador heredaba el
control indiviso de activos sin precio y ningún tribunal capaz de contener a nadie, y
la concentración resultante capturaba después el propio proceso electoral. Las
elecciones rusas de 1996 las ganó un titular financiado por los hombres a quienes
había vendido la economía; el problema fue la secuencia, no la papeleta
\citep{aslund2013}.

La propia historia cubana dice lo mismo desde el otro lado. El
capítulo~\ref{ch:inheritance} sostuvo que el rasgo decisivo de enero de 1959 fue que
el golpe de Batista ya había destruido toda institución capaz de arbitrar entre
contendientes, de modo que el único cuerpo con legitimidad era el movimiento
guerrillero y no había nada que frenara su consolidación. Una transición cubana que
produzca un actor legítimo y ningún otro terminará donde terminó 1959, cualesquiera
que sean sus intenciones. Las instituciones tienen que existir antes de la contienda,
o la contienda decide quién es su dueño.

Hay que decir dos cosas en contra de esta posición, porque es la posición bajo la
cual los gobiernos autoritarios prometen siempre elecciones para más adelante.

Primero, ``más adelante'' tiene que ser una fecha. Un argumento de secuencia sin
calendario publicado es un aplazamiento indefinido con una justificación, y el
gobierno cubano viene ofreciendo exactamente eso desde 1959. El dispositivo de
compromiso aquí es que el calendario queda anclado al principio, junto con las
instituciones, y es él mismo justiciable.

Segundo, cierta apertura política es \emph{precondición} de las instituciones y no
premio por ellas. Una judicatura no puede volverse independiente en secreto; un
organismo anticorrupción no puede funcionar sin una prensa que publique lo que
encuentra; un registro de la propiedad solo merece confianza si su contenido puede
disputarse en público. La secuencia no es entonces ley-y-después-política, sino un
orden específico \emph{dentro} de la apertura política: asociación, prensa,
información y tribunales temprano, porque las instituciones los necesitan;
elecciones nacionales competitivas en una fecha posterior fija, porque antes hay que
ponerle precio a los activos y hay que dejar que los partidos se construyan.

\begin{design}{El calendario político}
\label{d:calendar}
\begin{rules}
\rl{1}{Inmediato} Asociación legal, prensa y edición independientes, sindicatos
independientes, y el derecho a recibir y difundir información a través de las
fronteras. Son precondición de todas las instituciones que siguen y no le cuestan
nada al Estado conceder.

\rl{2}{Año uno} Reforma judicial (Diseño~\ref{d:courts}), un tribunal
contencioso-administrativo ante el cual un ciudadano o una empresa pueda demandar al
Estado, publicación de las cuentas del Estado, y una ley de partidos que fije las
condiciones de registro de las organizaciones políticas ---deliberadamente permisiva,
con financiamiento público y divulgación---.

\rl{3}{Años dos y tres} Elecciones municipales competitivas, celebradas primero y dos
veces, porque los municipios son donde los partidos aprenden a ser partidos, donde la
administración electoral aprende su oficio, y donde un error es sobrevivible.

\rl{4}{Año cinco} Elecciones competitivas a la asamblea nacional, en una fecha fijada
en la enmienda constitucional que inicia el proceso, con una comisión electoral
independiente constituida en el año uno y sin mayoría del oficialismo en ella.

\rl{5}{No aplazamiento} Las fechas quedan ancladas y su aplazamiento exige una
mayoría cualificada más la confirmación judicial de una emergencia declarada, con
caducidad automática. Un calendario que un gobierno puede mover a voluntad no es un
calendario.
\end{rules}
\end{design}

\section{Los tribunales}

Cuba no ha tenido desde 1959 un tribunal que fallara contra el ejecutivo en materia
política, y la práctica de subordinar la judicatura a la dirección política se
estableció en las primeras semanas de ese año, cuando se volvió a juzgar a acusados
absueltos por insistencia pública del líder. Sesenta y siete años de eso producen
una magistratura con competencia técnica en el trabajo civil y penal corriente y sin
ninguna memoria institucional de independencia.

La reforma no puede ser entonces un decreto que declare independientes a los jueces.
Tiene que cambiar las cuatro cosas que de verdad determinan si un juez puede fallar
contra el gobierno: quién los nombra, cuánto duran, quién les paga y quién puede
removerlos.

\begin{design}{Independencia judicial}
\label{d:courts}
\begin{rules}
\rl{1}{Nombramiento} Los jueces los nombra un consejo judicial de composición
mayoritariamente judicial, cuyos miembros no judiciales designa la asamblea por
mayoría cualificada y no el ejecutivo. Las vacantes se anuncian; los criterios y las
ternas se publican.

\rl{2}{Permanencia} Nombramiento hasta la edad de retiro, removible solo por
incapacidad o falta probada, mediante un procedimiento que es él mismo judicial y
público. Se suprime el arreglo actual, en el que los jueces sirven a voluntad de la
asamblea que los elige.

\rl{3}{Presupuesto} El presupuesto de la judicatura es una proporción mínima fija
del gasto estatal, fijada en la constitución, desembolsada bajo la autoridad de la
propia judicatura y auditada públicamente. Una judicatura cuya factura eléctrica
paga el ejecutivo a discreción no es independiente, y en Cuba esto no es una
metáfora.

\rl{4}{Tribunal administrativo} Se crea una jurisdicción contencioso-administrativa
especializada en la que cualquier persona o empresa pueda demandar a cualquier
órgano del Estado por una licencia denegada, un permiso demorado, un bien incautado,
un contrato incumplido o un reglamento mal aplicado, con el Estado asumiendo las
costas cuando pierde. Es el tribunal más importante de la Parte~III: es donde los
dispositivos de compromiso del Diseño~\ref{d:commit} se vuelven reales, y es el
mecanismo por el cual se hace cumplir la eliminación de la discrecionalidad del
capítulo~\ref{ch:integrity}.

\rl{5}{Publicación} Todas las sentencias se publican íntegras, en una base de datos
pública y gratuita, dentro de los treinta días. Un cuerpo de razonamiento publicado
es lo que convierte un conjunto de decisiones en un derecho, y es además el control
más barato disponible sobre las decisiones mismas.

\rl{6}{Capacidad mercantil} Las salas mercantiles, concursales y administrativas se
dotan con jueces formados para ellas ---un programa específico, financiado y de una
década, puesto que el capítulo~\ref{ch:socialstate} estableció que Cuba formó a muy
pocos juristas para pensar cómo funciona un mercado---. Es una restricción de
formación, no legislativa, y es vinculante.
\end{rules}
\end{design}

\section{Lo que tiene que ser legal}

Lo que sigue es un piso y no un ideal, enunciado como el mínimo del que depende el
resto de la Parte III. Cada punto está incluido porque algo posterior fracasa sin
él, y se da la razón.

\begin{design}{Las condiciones mínimas}
\label{d:minimum}
\begin{rules}
\rl{1}{Asociación} El derecho a formar organizaciones y afiliarse a ellas
---partidos, sindicatos, colegios profesionales, cooperativas, iglesias, grupos
vecinales--- mediante registro con criterios publicados, y con la denegación
apelable ante el tribunal administrativo. Lo exigen: todos los actores colectivos de
la Parte~III, incluidas las asociaciones de agricultores del
capítulo~\ref{ch:land} y los colegios profesionales que tienen que acreditar a los
reguladores del capítulo~\ref{ch:life}.

\rl{2}{Prensa} Edición y radiodifusión independientes, mediante registro, sin censura
previa y con la difamación como asunto civil y no penal. Lo exige: el
capítulo~\ref{ch:integrity}, cuyo modelo entero de aplicación es que alguien publique
lo que encuentre el auditor. Se derogan la Ley 88 de 1999 y las disposiciones penales
sobre propaganda enemiga.

\rl{3}{Información} Un derecho legal de acceso a la información en poder del Estado,
con presunción de divulgación, un plazo legal corto, una lista de excepciones
estrecha y enumerada, y apelación ante el tribunal administrativo. Lo exigen: la
contratación abierta, las declaraciones patrimoniales, el catastro publicado, el
registro electoral, y los datos ambientales del capítulo~\ref{ch:nature}.

\rl{4}{Movimiento} El derecho a salir y a regresar, anclado, con derogación de las
restricciones profesionales de salida que quedan. Lo exigen: la reforma de las
misiones médicas del capítulo~\ref{ch:life}, y el hecho simple de que un Estado que
puede impedir la salida nunca tiene que ganarse el consentimiento de sus ciudadanos
para quedarse.

\rl{5}{Comunicación} Ninguna interrupción del servicio de internet o telefonía por
motivos políticos, exigiéndose para cualquier interrupción autorización judicial,
publicación y duración declarada. Lo exige: el capítulo~\ref{ch:talent}, cuya
proposición entera es que un cliente extranjero pueda confiar en que un proveedor
cubano será localizable. El apagón nacional del 11 de julio de 2021 sería, bajo esta
regla, ilegal ---y ese es el sentido de escribirla---.
\end{rules}
\end{design}

\section{El trato con los actuales titulares}

El capítulo~\ref{ch:premises} estableció que el partido, las fuerzas armadas y GAESA
son parte de cualquier transición y no obstáculos que haya que retirar de ella. Esta
sección enuncia el trato, porque un trato que se deja implícito es uno que se
renegocia en cada paso.

El intercambio es: \emph{seguridad, posición y propiedad, a cambio de la entrega de
la discrecionalidad}. Los titulares conservan su libertad, sus pensiones, sus
viviendas, una vía legal hacia la política, y ---dentro de límites y por un proceso
transparente--- una participación en los activos que hoy manejan. Entregan el poder
de decidir casos, asignar licencias, fijar tipos de cambio, dirigir persecuciones
penales y controlar la divisa.

Esto no es generosidad y no debe presentarse como tal. Es el precio de una
transición que no pase por un derrumbe, y las historias alternativas están
disponibles: las transiciones que no ofrecieron nada a la élite saliente fueron las
violentas, y las que le ofrecieron todo produjeron oligarquías. El trabajo del
diseño es ser específico sobre dónde cae la línea.

\begin{design}{La amnistía, y dónde se detiene}
\label{d:amnesty}
\begin{rules}
\rl{1}{Cubierto} Amnistía por el hecho de haber ocupado un cargo, haber pertenecido
al partido, haber servido en los órganos de seguridad, haber administrado la
economía, haber aplicado las leyes entonces vigentes o haber participado en el
sistema político tal como existía. Nadie es procesado por haber formado parte de él.

\rl{2}{Cubierto} Amnistía por delitos económicos por debajo de un umbral publicado,
y por delitos derivados de la conducción ordinaria de una empresa estatal bajo las
reglas entonces aplicables.

\rl{3}{No cubierto} El homicidio, la tortura, la desaparición forzada, la violencia
sexual y las órdenes de cometerlos. No se amnistían, en ningún rango, y no corre
prescripción sobre ellos. Una transición que amnistía la tortura enseña que la
tortura es una opción de política pública.

\rl{4}{No cubierto} El enriquecimiento personal a gran escala por abuso de cargo, por
encima del umbral de la cláusula~2, cuando el activo pueda rastrearse. Aquí el
remedio es la recuperación del activo antes que la prisión, lo que es a la vez más
alcanzable y más útil.

\rl{5}{Condicionado} La amnistía de las cláusulas~1 y~2 queda condicionada a la
divulgación completa ante la comisión del Diseño~\ref{d:truth} y a la entrega de los
activos no declarados en el exterior. Es revocable por declaración materialmente
falsa, y solo por eso.

\rl{6}{Simétrico} La amnistía corre en las dos direcciones. Toda persona condenada
por expresión política, asociación, salida no autorizada o los hechos del 11 de
julio de 2021 es liberada y su condena anulada. Esto no es una concesión que se
canjee contra las demás; es una precondición, y ocurre primero.
\end{rules}
\end{design}

\begin{design}{Verdad sin persecución}
\label{d:truth}
\begin{rules}
\rl{1}{Mandato} Una comisión con un plazo fijo de tres años, el deber de recibir
testimonio de cualquier persona, potestades para requerir documentos al Estado, y el
deber de publicar. No tiene poder para procesar, para castigar, ni para señalar a
una persona como penalmente responsable.

\rl{2}{Alcance} Todo el período, en las dos direcciones, y los dos tipos de daño: las
ejecuciones y las condenas de prisión, los campos, las purgas, las confiscaciones y
las turbas de repudio; y también las muertes causadas por las organizaciones del
exilio, las campañas de sabotaje y Bahía de Cochinos. Una comisión que examina un
solo lado produce un documento que la mitad del país no leerá.

\rl{3}{Archivos} Los archivos del Estado sobre estas materias se transfieren a la
comisión y después a un archivo público, con un período de reserva fijo y breve por
privacidad personal genuina y por nada más.

\rl{4}{Restitución de la posición} La comisión puede recomendar, y el Estado debe
ejecutar, la restitución de la condición profesional, las pensiones y los
expedientes académicos a las personas apartadas por la parametración entre 1971 y
1976 y a sus equivalentes en períodos posteriores. Buena parte de ese daño se hizo
por carta administrativa y puede deshacerse igual, a costo casi nulo, y es el acto de
reparación más barato disponible.

\rl{5}{Sin jerarquía de víctimas} La comisión no clasifica sufrimientos, no arbitra
entre relatos ni emite un veredicto sobre la revolución. Produce un registro.
\end{rules}
\end{design}

\section{Dónde reside el poder}

Un punto estructural, fácil de omitir y caro de omitir.

El centralismo cubano es más viejo que la revolución. La república se gobernó desde
La Habana, la colonia se gobernó desde La Habana, y el Poder Popular de la
constitución de 1976 dio a los municipios elecciones sin darles dinero. Un municipio
que no puede recaudar ni retener ingresos no puede mantener una carretera, un
consultorio ni una bomba de agua, y así todos los problemas del país llegan al mismo
escritorio en la capital, donde hacen cola.

Esto importa para la Parte III específicamente porque mucho de lo que propone es
local: la generación y distribución en tejado del capítulo~\ref{ch:energy}, las
licencias y las compras de alimentos del capítulo~\ref{ch:tourism}, el catastro y el
control de la edificación del capítulo~\ref{ch:capital}, los retiros costeros del
capítulo~\ref{ch:nature}, y el impuesto sobre la propiedad del
capítulo~\ref{ch:welfare}, que es el ingreso local natural y el instrumento que hace
al gobierno local responsable ante los residentes locales en vez de ante el
ministerio.

\begin{design}{Capacidad municipal}
\label{d:municipal}
\begin{rules}
\rl{1}{Ingresos propios} Los municipios retienen el impuesto sobre la propiedad y una
proporción fija del gravamen al alojamiento y de las tasas de licencia de negocio
recaudadas en su territorio, fijada por fórmula y no por negociación.

\rl{2}{Igualación} Una fórmula de transferencia, publicada y basada en reglas,
redistribuye hacia las provincias orientales según indicadores medidos de necesidad.
Sin esto, la restricción regional del capítulo~\ref{ch:premises} se convierte en
divergencia regional: los ingresos del turismo se acumulan donde están los turistas,
y los turistas no están en Guantánamo.

\rl{3}{Competencias} Los municipios tienen y ejercen el control de la edificación, las
licencias locales, los viales locales, el agua y los residuos, y los locales de
atención primaria y de las escuelas, con los estándares de servicio fijados a nivel
nacional y la prestación a nivel local.

\rl{4}{Cuentas} Cada municipio publica su presupuesto, su ejecución y sus contratos,
en un formato común legible por máquina, trimestralmente. Es también el instrumento
más barato y más eficaz del capítulo~\ref{ch:integrity}, porque el robo local es
visible para la gente local si tiene los números disponibles.

\rl{5}{Movimiento interno} Se deroga el Decreto 217 de 1997 y las restricciones
asociadas a la residencia en La Habana. Restringir la migración interna no desarrolla
el oriente; solo convierte ser oriental en una discapacidad legal.
\end{rules}
\end{design}

\section{Lo que este capítulo no resuelve}

Dos cosas quedan abiertas, deliberadamente, y el lector debe saber cuáles.

La primera es la forma del Estado ---presidencial o parlamentario, unicameral o
bicameral, la fórmula electoral---. Importan, se estudian extensamente en otra
parte, y nada de la economía de este libro depende de ellas. Corresponden a quien
escriba la constitución y deberían resolverlas los cubanos en un proceso
constituyente, no proponerlas un economista en un capítulo.

La segunda es quién empieza. El diseño de este capítulo es ejecutable por el
gobierno actual ---cada una de sus cláusulas podría aprobarla la Asamblea Nacional
tal como está constituida, y varias de ellas, señaladamente la amnistía en las dos
direcciones y la liberación de las comunicaciones de la interferencia política, le
costarían muy poco y le comprarían mucho---. Es también ejecutable por un gobierno
sucesor. No es ejecutable por nadie, y el capítulo~\ref{ch:sequence} vuelve sobre la
cuestión de qué coalición lleva cada fase.
